\chapter{Generating Functions}

\begin{center}\itshape
In 1873 Francis Galton asked in the \emph{Educational Times} why so many family names of the English\\
aristocracy had died out. If each man has $0$, $1$ or $2$ sons with probabilities $\frac14$, $\frac14$, $\frac12$,\\
the population grows on average by a quarter per generation. Can the name still disappear?
\end{center}

A law on $\N$ is a sequence of numbers $p_0,p_1,p_2,\ldots$, and a sequence can be packed into a function:
$G(s)=\sum_kp_ks^k$. This chapter shows that the packing is worth it. Adding independent variables multiplies
their generating functions; a random number of random terms composes them; and the family tree of a population
in which everybody reproduces independently, the \emph{Galton--Watson process}, is described by iterating one
function. Galton's question then reduces to looking for the fixed points of a convex function on $[0,1]$.

The chapter ends with the Catalan numbers, which count Dyck paths, plane trees and the family trees of a
critical branching process, and with the one fact about power series that we need, proved from scratch.

\section{Probability generating functions}

\begin{definition}
Let $X$ be a random variable with values in $\N$ and $p_k=P(X=k)$. Its \emph{probability generating function}
(PGF) is
\[
G_X(s)=\mathbb E[s^X]=\sum_{k\ge0}p_ks^k,\qquad s\in[-1,1].
\]
\end{definition}

The series converges absolutely on $[-1,1]$ since $\sum p_k=1$, and the two expressions agree by the transfer
formula (Proposition~\ref{prop:transfer}). For instance, Bernoulli$(p)$ has PGF $1-p+ps$; binomial$(n,p)$ has
$(1-p+ps)^n$ by the binomial formula; geometric$(p)$ has $\sum_{k\ge1}(1-p)^{k-1}ps^k=\frac{ps}{1-(1-p)s}$; and
Poisson$(\lambda)$ has $\sum_ke^{-\lambda}\frac{(\lambda s)^k}{k!}=e^{\lambda(s-1)}$.

We need two facts about power series. Volume~1 did not prove that a power series can be differentiated term by term,
so we prove what we use.

\begin{lemma}\label{lem:powerseries}
Let $f(s)=\sum_{k\ge0}a_ks^k$ with $\sum_k|a_k|r^k<\infty$ for some $r>0$.
\begin{enumerate}[label=(\alph*)]
\item \emph{Uniqueness}: if $f(s)=0$ for all $s\in[0,\varepsilon[$ for some $\varepsilon>0$, then all $a_k=0$.
\item \emph{Differentiation}: on $]-r,r[$, $f$ is differentiable and $f'(s)=\sum_{k\ge1}ka_ks^{k-1}$, a series that satisfies the
same hypothesis with any $r'<r$ in place of $r$. So $f$ has derivatives of all orders on $]-r,r[$, obtained term by term.
\end{enumerate}
\end{lemma}
\begin{proof}
(a) Suppose not, and let $j$ be the smallest index with $a_j\ne0$. For $0<s<\min(\varepsilon,r)$,
$0=f(s)/s^j=a_j+\sum_{k>j}a_ks^{k-j}$, and $\bigl|\sum_{k>j}a_ks^{k-j}\bigr|\le\frac s{r^{j+1}}\sum_{k>j}|a_k|r^k$, which tends to $0$ as
$s\to0$. So $a_j=0$, a contradiction.

(b) Fix $\rho<r$. Since $k^2(\rho/r)^k\to0$, it is bounded by some $B$, so $\sum_kk^2|a_k|\rho^k\le B\sum_k|a_k|r^k<\infty$; in particular
the series $\sum ka_ks^{k-1}$ converges absolutely for $|s|\le\rho$. For $|s|,|t|\le\rho$, Taylor's formula with integral
remainder (Volume~1, Chapter~8) applied to $x\mapsto x^k$, whose second derivative is bounded by $k(k-1)\rho^{k-2}$ between $s$
and $t$, gives $|t^k-s^k-ks^{k-1}(t-s)|\le\frac{k(k-1)}2\rho^{k-2}|t-s|^2$. Summing,
\[
\begin{aligned}&\Bigl|f(t)-f(s)-(t-s)\sum_kka_ks^{k-1}\Bigr|\eqbr\le\frac{|t-s|^2}{2\rho^2}\sum_kk^2|a_k|\rho^k,\end{aligned}
\]
which is $o(|t-s|)$. This is the derivative formula at every $s$ with $|s|<\rho$, and $\rho<r$ is arbitrary.
\end{proof}

\begin{proposition}\label{prop:pgf}
Let $X$ take values in $\N$, with PGF $G$.
\begin{enumerate}[label=(\alph*)]
\item $G$ determines the law of $X$: $p_k=G^{(k)}(0)/k!$.
\item $G$ is continuous, non-decreasing and convex on $[0,1]$, with $G(1)=1$.
\item $\mathbb E[X]=\lim_{s\to1^-}G'(s)$ and $\mathbb E[X(X-1)]=\lim_{s\to1^-}G''(s)$, in $[0,+\infty]$.
\end{enumerate}
\end{proposition}
\begin{proof}
(a) By Lemma~\ref{lem:powerseries}(b) with $r=1$, $G^{(k)}(0)=k!\,p_k$.

(b) On $[0,1[$, $G'=\sum kp_ks^{k-1}\ge0$ and $G''=\sum k(k-1)p_ks^{k-2}\ge0$. At $s=1$: for $0\le s\le1$ and any $K$,
$|G(1)-G(s)|\le\sum_{k\le K}p_k(1-s^k)+\sum_{k>K}p_k$; the first sum tends to $0$ as $s\to1$ and the second is small for
$K$ large, so $G(s)\to G(1)=1$. Monotonicity and convexity extend to $[0,1]$ by continuity.

(c) Let $m=\sum_kkp_k\in[0,\infty]$. For $s<1$, $G'(s)\le m$. For each $K$, $G'(s)\ge\sum_{k\le K}kp_ks^{k-1}\to\sum_{k\le K}kp_k$ as $s\to1$,
so $\liminf_{s\to1}G'(s)\ge\sum_{k\le K}kp_k$ for every $K$, hence $\ge m$. Since $G'$ is non-decreasing, its limit at $1$ exists and
equals $m$. The same argument applies to $G''$ and $\sum k(k-1)p_k$.
\end{proof}

\section{Sums and random sums}

\begin{proposition}\label{prop:pgfsum}
\begin{enumerate}[label=(\alph*)]
\item If $X$ and $Y$ are independent with values in $\N$, then $G_{X+Y}=G_XG_Y$.
\item \emph{Random sums}: let $X_1,X_2,\ldots$ be independent with the same law, of PGF $G_X$, and $N$ be independent of
them, with values in $\N$ and PGF $G_N$. Then $S_N=X_1+\dots+X_N$ (with $S_0=0$) has PGF $G_N\circ G_X$.
\end{enumerate}
\end{proposition}
\begin{proof}
(a) For $s\in[-1,1]$, $s^X$ and $s^Y$ are bounded and independent (Proposition~\ref{prop:indepfun}), so
$\mathbb E[s^{X+Y}]=\mathbb E[s^X]\mathbb E[s^Y]$ by Theorem~\ref{thm:linearity}(c).

(b) By (a) and induction, $S_n=X_1+\dots+X_n$ has PGF $G_X^n$. By total probability over the values of $N$, and
because $S_n$ is a function of $X_1,\dots,X_n$, independent of $N$,
\[
\begin{aligned}P(S_N=k)&=\sum_nP(N=n,S_n=k)\eqbr=\sum_nP(N=n)P(S_n=k).\end{aligned}
\]
Multiply by $s^k$ with $0\le s\le1$ and sum over $k$; exchanging the sums of non-negative terms,
$G_{S_N}(s)=\sum_nP(N=n)G_X(s)^n=G_N(G_X(s))$. Two PGFs that agree on $[0,1]$ agree on $[-1,1]$ by
Proposition~\ref{prop:pgf}(a).
\end{proof}

\begin{corollary}\label{cor:wald}
If $X_1$ and $X_2$ are independent Poisson of parameters $\lambda,\mu$, then $X_1+X_2$ is Poisson$(\lambda+\mu)$. In the situation of
Proposition~\ref{prop:pgfsum}(b), \emph{Wald's identity} holds: $\mathbb E[S_N]=\mathbb E[N]\,\mathbb E[X_1]$.
\end{corollary}
\begin{proof}
$e^{\lambda(s-1)}e^{\mu(s-1)}=e^{(\lambda+\mu)(s-1)}$, and a PGF determines the law. For Wald, differentiate
$G_N(G_X(s))$ for $s<1$: $G_N'(G_X(s))G_X'(s)$, and let $s\to1$, using Proposition~\ref{prop:pgf}(c) and $G_X(s)\to1^-$ (if
$G_X(s)=1$ for some $s<1$, then $X=0$ almost surely and both sides vanish).
\end{proof}

\section{Galton--Watson processes}

\begin{definition}
Let $(\xi_{n,i})_{n\ge0,i\ge1}$ be independent random variables with the same law $(p_k)_{k\ge0}$ on $\N$, the
\emph{offspring law}, with PGF $G$ and mean $m=\sum kp_k$. The \emph{Galton--Watson process} is $Z_0=1$ and
$Z_{n+1}=\sum_{i=1}^{Z_n}\xi_{n,i}$: the $i$-th individual of generation $n$ has $\xi_{n,i}$ children. The event
$\{Z_n=0\text{ for some }n\}$ is \emph{extinction}.
\end{definition}

(The existence of infinitely many independent offspring variables is admitted, as in Chapter~1; every event
below involves finitely many of them.)

\begin{figure}[ht]
\centering
\fig{ch03-tree}
\caption{Seven generations of a Galton--Watson tree with offspring law $(\frac14,\frac14,\frac12)$ on $\{0,1,2\}$. Coral
individuals have no children. This tree survives, but each run of this process dies out with probability $\frac12$.}
\label{fig:gwtree}
\end{figure}

\begin{proposition}\label{prop:gwgen}
$Z_n$ has PGF $G_n=G\circ G\circ\dots\circ G$ ($n$ times), and $\mathbb E[Z_n]=m^n$ when $m<\infty$.
\end{proposition}
\begin{proof}
$Z_{n+1}$ is a random sum of the independent variables $\xi_{n,1},\xi_{n,2},\ldots$, with a random number of terms
$Z_n$ that is a function of the $\xi_{j,i}$ with $j<n$, hence independent of them. Proposition~\ref{prop:pgfsum}(b) gives
$G_{n+1}=G_n\circ G$, and by induction $G_n=G^{\circ n}$. Wald's identity gives $\mathbb E[Z_{n+1}]=m\,\mathbb E[Z_n]$.
\end{proof}

So the mean population grows geometrically when $m>1$. That does not prevent extinction: with probability $p_0$ the
first individual has no child at all.

\begin{think}
Let $q_n=P(Z_n=0)$. Express $q_{n+1}$ in terms of $q_n$ and $G$, then guess the limit of $q_n$ from a picture of
the graph of $G$ and the diagonal. What decides whether the graph crosses the diagonal before $s=1$?
\end{think}

\begin{theorem}[extinction]\label{thm:extinction}
The probability of extinction $q$ is the smallest solution of $G(s)=s$ in $[0,1]$. If $p_1<1$, then $q=1$ if
$m\le1$, and $q<1$ if $m>1$.
\end{theorem}
\begin{proof}
The events $\{Z_n=0\}$ increase with $n$ (an extinct population stays extinct), and their union is extinction, so
$q=\lim q_n$ with $q_n=P(Z_n=0)=G_n(0)$ by continuity (Proposition~\ref{prop:rules}(e)). Then
$q_{n+1}=G(G_n(0))=G(q_n)$, and letting $n\to\infty$ with $G$ continuous, $q=G(q)$. If $r\in[0,1]$ is any solution of
$G(r)=r$, then $q_0=0\le r$, and if $q_n\le r$ then $q_{n+1}=G(q_n)\le G(r)=r$ because $G$ is non-decreasing. So all
$q_n\le r$, and $q\le r$: $q$ is the smallest fixed point.

Let $f(s)=G(s)-s$ on $[0,1]$, so $f(1)=0$ and $f(0)=p_0\ge0$. If $p_0+p_1=1$, $G(s)=p_0+p_1s$ and $m=p_1<1$; then $f(s)=p_0(1-s)>0$ for
$s<1$ (as $p_0=1-p_1>0$), and $q=1$. Otherwise some $p_k>0$ with $k\ge2$, so $G''(s)>0$ on $]0,1[$ and $f'$ is increasing
there, with $f'(s)\to m-1$ as $s\to1$ (Proposition~\ref{prop:pgf}(c)).

If $m\le1$, then $f'(s)<m-1\le0$ on $]0,1[$, so $f$ is decreasing on $[0,1]$ and $f(s)>f(1)=0$ for $s<1$: the only fixed point
is $1$, and $q=1$. If $m>1$ (possibly $m=\infty$), then $f'(s)>0$ for $s$ close to $1$, so $f(s)<f(1)=0$ there; since $f(0)\ge0$,
the intermediate value theorem gives a fixed point in $[0,1[$, and $q<1$.
\end{proof}

\begin{figure}[ht]
\centering
\fig{ch03-cobweb}
\caption{The iteration $q_{n+1}=G(q_n)$ from $q_0=0$ climbs the staircase between the graph of $G$ and the diagonal.
Left: offspring law $(\frac12,\frac14,\frac14)$, mean $\frac34$, and the staircase goes to $1$. Right: $(\frac14,\frac14,\frac12)$, mean $\frac54$,
and it stops at the first crossing $q=\frac12$.}
\label{fig:cobweb}
\end{figure}

For Galton's question, $G(s)=\frac14+\frac14s+\frac12s^2$, and $G(s)=s$ becomes $2s^2-3s+1=0$, with roots $\frac12$ and $1$: despite
an average growth of $25\%$ per generation, the name dies out with probability $\frac12$ (Figure~\ref{fig:cobweb}). In the
\emph{critical} case $m=1$ extinction is certain even though the mean population stays equal to $1$: rare runs
grow large and compensate the many that die. In 1931, Alfred Lotka fitted the law
$p_0=0.4825$, $p_k=0.2126\times0.5893^{k-1}$ ($k\ge1$) to American census data on sons; its mean is about $1.26$ and the
extinction probability of a male line is about $0.82$ (Exercise~\ref{ex:lotka}).

\begin{remark}
On survival, the population does not linger: Chapter~10 proves that $Z_n/m^n$ converges, and that for $m>1$ and a
finite offspring variance its limit is positive exactly on survival; the Kesten--Stigum theorem, quoted there, weakens the
variance condition to $\mathbb E[\xi\ln\xi]<\infty$. So a surviving population grows like $m^n$, as the naive mean suggests.
\end{remark}

\section{Catalan numbers, paths and trees}

The Catalan numbers $C_n=\frac1{n+1}\binom{2n}n$ count the Dyck paths of length $2n$ (Exercise~\ref{ex:catalan}). Generating
functions give them a second life.

\begin{proposition}\label{prop:catalanrec}
$C_0=1$ and $C_{n+1}=\sum_{k=0}^nC_kC_{n-k}$. Consequently $C(x)=\sum_nC_nx^n$ satisfies $C(x)=1+xC(x)^2$ for
$|x|\le\frac14$, and $C(\frac14)=2$.
\end{proposition}
\begin{proof}
Decompose a Dyck path of length $2n+2$ at its first return to $0$, at time $2k+2$: it is an up-step, a Dyck path of length
$2k$ lifted by one, a down-step, then a Dyck path of length $2(n-k)$. This is a bijection with pairs of Dyck paths of
lengths $2k$ and $2(n-k)$, which gives the recursion. By Stirling's formula (Volume~1, Theorem~3.17),
$C_n\sim\frac{4^n}{\sqrt\pi n^{3/2}}$, so $\sum C_n4^{-n}<\infty$ and $C$ converges absolutely for $|x|\le\frac14$. The recursion says that the
coefficients of $xC(x)^2$ (computed by the Cauchy product of absolutely convergent series, which is a sum by packets)
are those of $C(x)-1$. At $x=\frac14$, $c=C(\frac14)$ satisfies $c^2-4c+4=0$, so $c=2$.
\end{proof}

Solving the quadratic gives $C(x)=\frac{1-\sqrt{1-4x}}{2x}$, the sign being fixed by $C(0)=1$; Exercise~\ref{ex:binseries}
extracts $C_n$ from it again.

A \emph{plane tree} is a rooted tree in which the children of each vertex are ordered from left to right. Walk around
a plane tree with $n$ edges, starting at the root, going down the leftmost unexplored edge whenever possible and up
otherwise, and record the depth after each step: this \emph{contour walk} crosses every edge once down and once up,
and never goes below the root. It is a Dyck path of length $2n$ (Figure~\ref{fig:contour}).

\begin{figure}[ht]
\centering
\fig{ch03-contour}
\caption{A plane tree with $9$ vertices and its contour walk, a Dyck path of length $16$.}
\label{fig:contour}
\end{figure}

\begin{proposition}\label{prop:planetrees}
The contour walk is a bijection between plane trees with $n$ edges and Dyck paths of length $2n$. So there are
$C_n$ plane trees with $n+1$ vertices.
\end{proposition}
\begin{proof}
By induction on $n$, the case $n=0$ being clear. A plane tree with $n\ge1$ edges is a root whose first child carries a
plane tree $A$, together with the plane tree $B$ formed by the root and its other children's subtrees. Its contour walk
is: up, the contour walk of $A$ lifted by one, down, then the contour walk of $B$. This is exactly the first-return
decomposition of Proposition~\ref{prop:catalanrec}, and it is reversible: from a Dyck path, the first return determines the
pair of shorter Dyck paths, hence by induction the pair of trees $(A,B)$, hence the tree.
\end{proof}

Now the link with branching. Take the offspring law $p_k=2^{-(k+1)}$, $k\ge0$ (geometric on $\N$, mean $1$, PGF $\frac1{2-s}$).
Since $m=1$, the tree is finite with probability $1$, and it is a random plane tree (order the children of each individual).

\begin{theorem}\label{thm:criticalgw}
For the Galton--Watson process with offspring law $p_k=2^{-(k+1)}$, every plane tree with $n$ vertices has probability
$2^{-(2n-1)}$ of being the family tree. Hence the total progeny $T$ satisfies $P(T=n)=C_{n-1}\,2^{-(2n-1)}$.
\end{theorem}
\begin{proof}
The family tree equals a given plane tree $t$ iff every vertex $v$ of $t$ has exactly $k_v$ children, its number of
children in $t$; by independence this has probability $\prod_v2^{-(k_v+1)}=2^{-\sum_v(k_v+1)}$. A tree with $n$ vertices has
$n-1$ edges, each counted once as a child, so $\sum_vk_v=n-1$ and the exponent is $2n-1$. Summing over the $C_{n-1}$ plane trees
with $n$ vertices gives the law of $T$.
\end{proof}

Since extinction is certain, these probabilities add up to $1$, which is the identity $\frac12C(\frac14)=1$ again. And since
\[
P(T=n)\sim\frac1{2\sqrt\pi}\,n^{-3/2},
\]
the mean total progeny is infinite: a critical population dies out surely, but after a
total number of individuals with infinite mean.

\begin{beauty}
One quadratic equation, $C=1+xC^2$, describes Dyck paths, plane trees, the triangulations of a polygon, the ways to
bracket a product, the stack-sortable permutations and the family trees of a critical population, and Stanley's
book lists more than two hundred families counted by the same numbers. Generating functions see the common
structure that the pictures hide: an object is either empty, or a root with two smaller objects of the same kind.
\end{beauty}

\section{Exercises}

\begin{exercise}[$\star$]\label{ex:thinning}
\emph{Thinning.} A Poisson$(\lambda)$ number of eggs are laid, and each hatches independently with probability $p$. Using
Proposition~\ref{prop:pgfsum}(b), show that the number of hatched eggs is Poisson$(\lambda p)$.
\end{exercise}

\begin{exercise}[$\star$]\label{ex:golden}
Each individual has $0$ or $3$ children with probability $\frac12$ each. Show that the extinction probability is
$\frac{\sqrt5-1}2=\frac1\varphi$.
\end{exercise}

\begin{exercise}[$\star$]\label{ex:lotka}
For Lotka's offspring law $p_0=0.4825$, $p_k=0.2126\times0.5893^{k-1}$ ($k\ge1$), compute $G$ in closed form and find the
extinction probability. (The published numbers are rounded, so $\sum p_k=1.00015$; ignore this.)
\end{exercise}

\begin{exercise}[$\star\star$]\label{ex:randomsumvar}
In the situation of Proposition~\ref{prop:pgfsum}(b), assume $\mathbb E[N^2]$ and $\mathbb E[X^2]$ finite. Show that
$\operatorname{Var}S_N=\mathbb E[N]\operatorname{Var}X+\operatorname{Var}N\,(\mathbb EX)^2$, and deduce $\operatorname{Var}Z_n$ for a Galton--Watson process with $m\ne1$ and offspring
variance $\sigma^2$.
\end{exercise}

\begin{exercise}[$\star\star$]\label{ex:progeny}
Let $T=\sum_{n\ge0}Z_n$ be the total progeny of a Galton--Watson process with $m<1$. Show that $\mathbb E[T]=\frac1{1-m}$, and that
its PGF $H$ satisfies $H(s)=s\,G(H(s))$ for $s\in[0,1]$.
\end{exercise}

\begin{exercise}[$\star\star$]\label{ex:firstpassage}
A walk starts at $0$ and moves $+1$ with probability $p$ and $-1$ with probability $q=1-p$, independently. Let $\tau$ be the
first time it reaches $1$ ($\tau=\infty$ if never). Show that $F(s)=\mathbb E[s^\tau\mathbf 1_{\tau<\infty}]$ satisfies $F(s)=ps+qsF(s)^2$ for
$0\le s<1$, deduce $F$, and show that $P(\tau<\infty)=\min(1,\frac pq)$.
\end{exercise}

\begin{exercise}[$\star\star$]\label{ex:binseries}
(a) For $|x|<\frac14$, let $g(x)=\sum_{n\ge0}\binom{1/2}n(-4x)^n$, where $\binom{1/2}n=\frac{\frac12(\frac12-1)\cdots(\frac12-n+1)}{n!}$. Show that
$(1-4x)g'(x)=-2g(x)$ and deduce $g(x)=\sqrt{1-4x}$. (b) Deduce from $C(x)=\frac{1-\sqrt{1-4x}}{2x}$ that $C_n=\frac1{n+1}\binom{2n}n$.
\end{exercise}

\begin{exercise}[$\star\star\star$]\label{ex:poissongw}
Let the offspring law be Poisson$(c)$ with $c>1$, and $\rho=1-q$ the survival probability. (a) Show that $\rho$ is the unique
solution in $]0,1[$ of $1-\rho=e^{-c\rho}$, and compute it for $c=2$. (b) Show that $\rho\sim2(c-1)$ as $c\to1^+$. (Chapter~12 meets this
number again as the size of the giant component of a random graph.)
\end{exercise}

\begin{exercise}[$\star\star$]\label{pb:sicherman}
Find all pairs of dice with six faces each, every face carrying a positive integer (repetitions allowed), such that the sum
of the two dice, thrown independently, has the same law as the sum of two ordinary dice.
\end{exercise}

\begin{exercise}[$\star\star$]\label{pb:loaded}
Two dice have faces $1,2,\dots,n$, with $n\ge2$, but they are loaded: the first shows $i$ with probability $p_i$, the
second $j$ with probability $q_j$, independently. Show that the sum cannot be uniformly distributed on $\{2,3,\dots,2n\}$.
Give two proofs for $n=6$: one with the roots of a polynomial, one with an inequality.
\end{exercise}

\begin{exercise}[$\star\star$]\label{pb:subsetsum}
\begin{enumerate}[label=(\alph*)]
\item Let $p$ be an odd prime. How many subsets of $\{1,2,\dots,p\}$ (the empty set included) have a sum divisible by $p$?
\item How many $p$-element subsets of $\{1,2,\dots,2p\}$ have a sum divisible by $p$?
\item Let $n\ge1$ and $0\le k\le n$, and let $N(n,k)$ be the number of $k$-element subsets of $\Z/n\Z$ whose sum is $0$.
Show that
\[
N(n,k)=\frac1n\sum_{d\mid\gcd(n,k)}(-1)^{k+k/d}\,\varphi(d)\binom{n/d}{k/d}.
\]
\end{enumerate}
\end{exercise}

\section{Hints}

\begin{hint}{ex:thinning}
The PGF of Bernoulli$(p)$ is $1-p+ps$; compose.
\end{hint}
\begin{hint}{ex:golden}
$1$ is always a root of $G(s)-s$; factor it out.
\end{hint}
\begin{hint}{ex:lotka}
$G(s)=p_0+\frac{as}{1-rs}$; the fixed-point equation is quadratic.
\end{hint}
\begin{hint}{ex:randomsumvar}
Differentiate $G_N\circ G_X$ twice and use $\mathbb E[X(X-1)]=G''(1^-)$.
\end{hint}
\begin{hint}{ex:progeny}
Condition on the number of children of the root: each child starts an independent copy of the whole tree.
\end{hint}
\begin{hint}{ex:firstpassage}
To go from $0$ to $1$, either step up at once, or step down and then go up twice, each time like a fresh walk.
\end{hint}
\begin{hint}{ex:binseries}
(a) Use Lemma~\ref{lem:powerseries}(b) and check the recursion of the coefficients; then differentiate $g(x)/\sqrt{1-4x}$.
\end{hint}
\begin{hint}{ex:poissongw}
(a) $G(s)=e^{c(s-1)}$. (b) Expand $e^{-c\rho}$ to second order in $\rho$.
\end{hint}

\begin{hint}{pb:sicherman}
Encode a die by the polynomial $\sum_{\text{faces}}x^{\text{face}}$; the sum of independent dice corresponds to the product.
Factor $(x+x^2+\cdots+x^6)^2$ into irreducibles of $\Q[X]$, and evaluate the factors at $0$ and $1$.
\end{hint}

\begin{hint}{pb:loaded}
For $n=6$: the generating polynomials have degree $5$, and $1+x+\dots+x^{10}$ has no real root. For every $n$: compare the
coefficients of $x^0$, $x^{2n-2}$ and $x^{n-1}$ of the product, with the AM--GM inequality.
\end{hint}

\begin{hint}{pb:subsetsum}
Expand $\prod_{a\in\Z/n\Z}(1+y\,\omega^{ja})$ and filter with $\frac1n\sum_j$. For a root of unity $\eta$ of order $d$,
$\prod_{\varepsilon^d=1}(1+y\varepsilon)=1-(-y)^d$.
\end{hint}

\section{Solutions}

\begin{solution}{ex:thinning}
\idea A thinned population is a random sum of Bernoulli variables, so its PGF is a composition.
\writeup The number of hatched eggs is $S_N=X_1+\dots+X_N$ with $N$ Poisson$(\lambda)$ and $X_i$ Bernoulli$(p)$, independent. Its PGF
is $G_N(G_X(s))=e^{\lambda(1-p+ps-1)}=e^{\lambda p(s-1)}$, the PGF of Poisson$(\lambda p)$, and a PGF determines the law. (Moreover the
hatched and unhatched counts are independent Poisson variables: see Chapter~9.)
\end{solution}

\begin{solution}{ex:golden}
\idea Write the fixed-point equation and remove the trivial root $s=1$.
\writeup $G(s)=\frac12+\frac12s^3$, and $G(s)=s$ iff $s^3-2s+1=0$, i.e.\ $(s-1)(s^2+s-1)=0$. The roots of $s^2+s-1$ are $\frac{-1\pm\sqrt5}2$,
and the one in $[0,1]$ is $\frac{\sqrt5-1}2=\frac1\varphi\approx0.618$. Since $m=\frac32>1$, Theorem~\ref{thm:extinction} says it is the extinction
probability.
\end{solution}

\begin{solution}{ex:lotka}
\idea The geometric tail sums to a rational function, so the fixed-point equation is a quadratic.
\writeup With $a=0.2126$ and $r=0.5893$, $G(s)=p_0+\sum_{k\ge1}ar^{k-1}s^k=p_0+\frac{as}{1-rs}$. Then $G(s)=s$ becomes
$(p_0-s)(1-rs)+as=0$, i.e.\ $rs^2-(1+rp_0-a)s+p_0=0$. Its roots are about $0.8193$ and $0.9994$ (the latter is $1$ up to the rounding of
the data), so $q\approx0.82$: more than four male lines out of five eventually die out, although $m=\frac a{(1-r)^2}\approx1.26$.
\end{solution}

\begin{solution}{ex:randomsumvar}
\idea Variances come from second derivatives of PGFs at $1$.
\writeup Write $H=G_N\circ G_X$. For $s<1$, $H''=G_N''(G_X)G_X'^2+G_N'(G_X)G_X''$. Letting $s\to1$ (Proposition~\ref{prop:pgf}(c)),
$\mathbb E[S_N(S_N-1)]=\mathbb E[N(N-1)]\mu^2+\mathbb E[N]\,\mathbb E[X(X-1)]$ with $\mu=\mathbb EX$. Adding $\mathbb E[S_N]=\mathbb E[N]\mu$ and subtracting $(\mathbb E[N]\mu)^2$,
\[
\begin{aligned}\operatorname{Var}S_N&=\mu^2\bigl(\mathbb E[N^2]-\mathbb E[N]^2\bigr)\eqbr+\mathbb E[N]\bigl(\mathbb E[X^2]-\mu^2\bigr)\\&=\mu^2\operatorname{Var}N+\mathbb E[N]\operatorname{Var}X.\end{aligned}
\]
For Galton--Watson, $v_{n+1}=m^2v_n+m^n\sigma^2$ with $v_0=0$. Dividing by $m^{2(n+1)}$, $\frac{v_{n+1}}{m^{2n+2}}=\frac{v_n}{m^{2n}}+\frac{\sigma^2}{m^{n+2}}$, so
$v_n=\sigma^2m^{2n-2}\sum_{j=0}^{n-1}m^{-j}=\sigma^2m^{n-1}\frac{m^n-1}{m-1}$.
\end{solution}

\begin{solution}{ex:progeny}
\idea The tree below each child of the root is a copy of the whole tree, independent of the others.
\writeup $\mathbb E[T]=\sum_n\mathbb E[Z_n]=\sum_nm^n=\frac1{1-m}$ (Proposition~\ref{prop:gwgen} and linearity for non-negative terms). For the PGF,
$T=1+T_1+\dots+T_{\xi}$, where $\xi=\xi_{0,1}$ is the number of children of the root and $T_i$ is the total progeny of the $i$-th
child. The $T_i$ are independent copies of $T$, independent of $\xi$: the subtree of child $i$ is built from offspring
variables that no other subtree uses. Here the variables $\xi_{n,i}$ are re-indexed along the tree, which does not change
their joint law; $T$ is finite almost surely since $q=1$. By Proposition~\ref{prop:pgfsum}(b), $\mathbb E[s^{T}]=s\,G(H(s))$.
\end{solution}

\begin{solution}{ex:firstpassage}
\idea First-step analysis: after one step down, reaching $1$ means reaching $0$ and then $1$, two independent copies of
the original problem.
\writeup With probability $p$ the first step reaches $1$: $\tau=1$. With probability $q$ it goes to $-1$; then $\tau=1+\tau'+\tau''$, where $\tau'$
is the time to climb from $-1$ to $0$ and $\tau''$ from $0$ to $1$. By the independence of the steps, $\tau'$ and $\tau''$ are independent
copies of $\tau$ (each depends on its own block of steps, and the walk is the same from any starting point). So, for
$0\le s<1$, $F(s)=ps+qsF(s)^2$. Solving, $F(s)=\frac{1-\sqrt{1-4pqs^2}}{2qs}$, the root that stays bounded as $s\to0$ (the other one
tends to $+\infty$). Letting $s\to1$, $P(\tau<\infty)=F(1^-)=\frac{1-\sqrt{1-4pq}}{2q}=\frac{1-|p-q|}{2q}$, since
$1-4pq=(p+q)^2-4pq=(p-q)^2$. If $p\ge q$ this is $\frac{2q}{2q}=1$; if $p<q$ it is $\frac{2p}{2q}=\frac pq$.
\end{solution}

\begin{solution}{ex:binseries}
\idea A power series is pinned down by a first-order linear differential equation, and the square root solves the
same one.
\writeup (a) Let $b_n=\binom{1/2}n(-4)^n$. Then $\frac{b_{n+1}}{b_n}=\frac{(\frac12-n)(-4)}{n+1}=\frac{4n-2}{n+1}$, so $|b_n|\le4^n$ and $g$ converges for $|x|<\frac14$.
By Lemma~\ref{lem:powerseries}(b), $(1-4x)g'(x)+2g(x)=\sum_n\bigl((n+1)b_{n+1}-4nb_n+2b_n\bigr)x^n=0$, since $(n+1)b_{n+1}=(4n-2)b_n$. Then
$h(x)=g(x)(1-4x)^{-1/2}$ has $h'=(1-4x)^{-3/2}\bigl((1-4x)g'+2g\bigr)=0$, so $h=h(0)=1$.
(b) $1-\sqrt{1-4x}=-\sum_{n\ge1}b_nx^n$, so by uniqueness of coefficients $C_n=-\frac{b_{n+1}}2$. Now
$-b_{n+1}=-\binom{1/2}{n+1}(-4)^{n+1}=\frac{2^{n+1}\cdot1\cdot3\cdots(2n-1)}{(n+1)!}$, using $\frac12(\frac12-1)\cdots(\frac12-n)=(-1)^n\frac{1\cdot3\cdots(2n-1)}{2^{n+1}}$. Since
$1\cdot3\cdots(2n-1)=\frac{(2n)!}{2^nn!}$, this gives $C_n=\frac{(2n)!}{(n+1)!\,n!}=\frac1{n+1}\binom{2n}n$.
\end{solution}

\begin{solution}{ex:poissongw}
\idea The survival probability is the non-trivial crossing of $e^{c(s-1)}$ with the diagonal; near criticality the crossing is
close to $1$ and a second-order expansion finds it.
\writeup (a) $G(s)=\sum e^{-c}\frac{c^ks^k}{k!}=e^{c(s-1)}$ and $m=c>1$, so $q<1$ is the smallest fixed point, and by the proof of
Theorem~\ref{thm:extinction} it is the only one in $[0,1[$ ($f=G-\mathrm{id}$ is strictly convex, so it has at most two zeros). With $\rho=1-q$:
$1-\rho=e^{-c\rho}$. For $c=2$, iterating $\rho\mapsto1-e^{-2\rho}$ from $\rho=1$ gives $\rho\approx0.7968$.
(b) Let $\varphi(\rho)=1-\rho-e^{-c\rho}$. Since $e^{-u}=1-u+\frac{u^2}2-\frac{u^3}6e^{-\theta u}$ for some $\theta\in]0,1[$, $\varphi(\rho)=(c-1)\rho-\frac{c^2\rho^2}2+O(\rho^3)$ with a
uniform constant for $c\le2$. As $c\to1^+$, $\rho\to0$ (otherwise along a subsequence $\rho\to\rho_0>0$ with $1-\rho_0=e^{-\rho_0}$, impossible
since $e^{-u}>1-u$ for $u>0$). Dividing $\varphi(\rho)=0$ by $\rho$, $(c-1)=\frac{c^2\rho}2+O(\rho^2)$, so $\rho\sim\frac{2(c-1)}{c^2}\sim2(c-1)$.
\end{solution}

\begin{solution}{pb:sicherman}
\idea A die is a multiset of faces, and the law of a sum of independent variables is the product of the generating functions
(Volume~4, Chapter~3). So the question is how $(x+\dots+x^6)^2$ factors. Unique factorisation in $\Q[X]$ leaves only finitely
many candidates, and the conditions ``six faces'' and ``positive faces'' are the values at $1$ and at $0$.

\writeup To a die with faces $a_1,\dots,a_6$ attach $A(x)=\sum_ix^{a_i}$; its coefficients are non-negative integers,
$A(1)=6$ and, since the faces are positive, $A(0)=0$. If $B$ is the polynomial of the second die, the number of ways to obtain
the sum $s$ is the coefficient of $x^s$ in $A(x)B(x)$, so the condition is
\[
A(x)B(x)=(x+x^2+\dots+x^6)^2=x^2(1+x)^2(1+x+x^2)^2(1-x+x^2)^2,
\]
because $x+\dots+x^6=x\frac{x^6-1}{x-1}=x(1+x)(1+x+x^2)(1-x+x^2)$. The four factors $x$, $1+x$, $1+x+x^2$, $1-x+x^2$ are
irreducible in $\Q[X]$: the quadratics have negative discriminant, hence no rational root. Since $\Q[X]$ has unique
factorisation (Volume~1, Chapter~6), $A=\lambda x^a(1+x)^b(1+x+x^2)^c(1-x+x^2)^d$ with $\lambda\in\Q^*$ and $0\le a,b,c,d\le2$,
and $B=\lambda'x^{2-a}(1+x)^{2-b}(1+x+x^2)^{2-c}(1-x+x^2)^{2-d}$ with $\lambda\lambda'=1$. The leading coefficients $\lambda$
and $\lambda'$ of $A$ and $B$ are positive integers, so $\lambda=\lambda'=1$.

Now evaluate. At $x=1$ the factors take the values $1,2,3,1$, so $A(1)=2^b3^c=6$ forces $b=c=1$. At $x=0$, $A(0)=0$ forces
$a\ge1$ and $B(0)=0$ forces $2-a\ge1$, so $a=1$. Only $d\in\{0,1,2\}$ is free. For $d=1$ both dice are ordinary. For $d=0$,
$A=x(1+x)(1+x+x^2)=x+2x^2+2x^3+x^4$, the die $1$, $2$, $2$, $3$, $3$, $4$, and then $d'=2$ gives
$B=x(1+x)(1+x+x^2)(1-x+x^2)^2=x+x^3+x^4+x^5+x^6+x^8$, the die $1$, $3$, $4$, $5$, $6$, $8$; $d=2$ is the same pair in the other order. Both
polynomials have non-negative coefficients, so they are dice. There are exactly two solutions: the ordinary pair and
Sicherman's pair.
\tools Volume~4, Chapter~3 (generating functions of sums); Volume~1, Chapter~6 (unique factorisation in $\Q[X]$).
\leads The same method classifies ``Sicherman dice'' with $n$ faces: there are as many pairs as ways of splitting the
cyclotomic factors $\Phi_d$, $d\mid n$, $d>1$, of $x\frac{x^n-1}{x-1}$ between two dice so that both have non-negative
coefficients (Gallian and Rusin, 1979).
\end{solution}

\begin{solution}{pb:loaded}
\idea Again the generating polynomials multiply. A real polynomial of odd degree has a real root, while the polynomial of the
uniform law on $2,\dots,12$ has none: that settles $n=6$. For general $n$, only three coefficients matter: the extreme ones are
products $p_1q_1$ and $p_nq_n$, and the middle one contains $p_1q_n+p_nq_1$, which AM--GM bounds from below.

\writeup Let $A(x)=\sum_{i=1}^np_ix^{i-1}$ and $B(x)=\sum_jq_jx^{j-1}$. If the sum were uniform on $\{2,\dots,2n\}$, then
\[
A(x)B(x)=\frac1{2n-1}\bigl(1+x+\dots+x^{2n-2}\bigr).
\]
\emph{First proof, $n=6$.} The right-hand side has degree $10$, so $A$ and $B$, of degree at most $5$, both have degree $5$.
A real polynomial of odd degree has a real root (intermediate value theorem), so $AB$ would have a real root. But
$1+x+\dots+x^{10}=\frac{x^{11}-1}{x-1}$ for $x\ne1$, whose roots are the $11$-th roots of unity other than $1$, none of them
real, and its value at $1$ is $11$.

\emph{Second proof, every $n\ge2$.} Comparing the coefficients of $x^0$ and $x^{2n-2}$ gives $p_1q_1=p_nq_n=\frac1{2n-1}$; in
particular all four numbers are positive. The coefficient of $x^{n-1}$ is $\sum_{i+j=n+1}p_iq_j\ge p_1q_n+p_nq_1$, and by
AM--GM
\[
p_1q_n+p_nq_1\ge2\sqrt{p_1q_np_nq_1}=2\sqrt{p_1q_1\cdot p_nq_n}=\frac2{2n-1}>\frac1{2n-1},
\]
a contradiction with the coefficient $\frac1{2n-1}$ of $x^{n-1}$ on the right.
\tools Volume~4, Chapter~3 (generating functions); Volume~1, Chapter~9 (AM--GM) and Chapter~3 (roots of polynomials).
\end{solution}

\begin{solution}{pb:subsetsum}
\idea ``The sum is divisible by $n$'' is a filter: average $\omega^{j\cdot\text{sum}}$ over $j$. The generating function
of subsets is a product, $\prod_a(1+y\,x^a)$, and at $x=\omega^j$ the product runs over the powers of a root of unity of
order $d$, each repeated $\frac nd$ times. Such products collapse to $1-(-y)^d$.

\writeup Fix $n$ and let $\omega=e^{2i\pi/n}$. For a subset $A$ of $\Z/n\Z$, $\frac1n\sum_{j=0}^{n-1}\omega^{j\sum A}$ is $1$
if $\sum A=0$ and $0$ otherwise (Volume~1, Chapter~5). So
\[
N(n,k)=\frac1n\sum_{j=0}^{n-1}[y^k]F_j(y),\qquad F_j(y)=\prod_{a\in\Z/n\Z}(1+y\,\omega^{ja}),
\]
because expanding $F_j$ gives $\sum_Ay^{|A|}\omega^{j\sum A}$. Let $d$ be the order of $\omega^j$, that is,
$d=n/\gcd(n,j)$. As $a$ runs over $\Z/n\Z$, $\omega^{ja}$ runs over the $d$-th roots of unity, each $\frac nd$ times. Now
$\prod_{\varepsilon^d=1}(X-\varepsilon)=X^d-1$; put $X=-\frac1y$ and multiply by $(-y)^d$:
$\prod_\varepsilon(1+y\varepsilon)=1-(-y)^d$. Hence
\[
F_j(y)=\bigl(1-(-y)^d\bigr)^{n/d}=\sum_{i}\binom{n/d}i(-1)^i(-1)^{di}y^{di}.
\]
The coefficient of $y^k$ is $0$ unless $d\mid k$, and then it is $(-1)^{k/d+k}\binom{n/d}{k/d}$. Exactly $\varphi(d)$ values
of $j$ give the order $d$, for each $d\mid n$. Summing gives the formula of (c).

For (a), $\{1,\dots,p\}$ is $\Z/p\Z$, and we count subsets of all sizes, so we take $y=1$: $F_0(1)=2^p$ and, for $j\ne0$,
$d=p$ and $F_j(1)=1-(-1)^p=2$. The number is $\frac1p\bigl(2^p+2(p-1)\bigr)=\frac{2^p-2}p+2$. For $p=3$ this is $4$:
$\varnothing$, $\{3\}$, $\{1,2\}$, $\{1,2,3\}$.

For (b), $\{1,\dots,2p\}$ contains each residue twice, so the count is $\frac1p\sum_j[y^p]F_j(y)^2$. For $j=0$ this is
$[y^p](1+y)^{2p}=\binom{2p}p$; for $j\ne0$ it is $[y^p](1+y^p)^2=2$. The answer is $\frac1p\bigl(\binom{2p}p+2(p-1)\bigr)
=\frac{\binom{2p}p-2}p+2$. For $p=5$ this is $52$.
\tools Volume~1, Chapter~5 (roots of unity); Volume~4, Chapter~3 (generating functions).
\leads Part (b) is Problem~6 of the 1995 International Mathematical Olympiad, where most solutions were long counting
arguments. The formula of (c) also counts binary necklaces with a given number of black beads and a ``sum'' condition; the
same filter computes the number of solutions of $a_1+\dots+a_k=0$ in any finite abelian group, and it is the starting point
of the circle method of Hardy and Littlewood (Volume~14).
\end{solution}

\section*{Sources and further reading}

\begin{itemize}
\item H.~W.~Watson and F.~Galton, \emph{On the probability of the extinction of families}, Journal of the
Anthropological Institute 4 (1875), 138--144. The original paper, whose conclusion that every family dies out was wrong
for $m>1$; the correct statement is due to J.~F.~Steffensen (1930).
\item A.~J.~Lotka, \emph{The extinction of families}, Journal of the Washington Academy of Sciences 21 (1931),
377--380 and 453--459.
\item K.~B.~Athreya and P.~E.~Ney, \emph{Branching Processes}, Springer, 1972. The standard reference, including the
Kesten--Stigum theorem.
\item H.~S.~Wilf, \emph{generatingfunctionology}, 3rd ed., A K Peters, 2006 (free online). Generating functions as a
way of life.
\item R.~P.~Stanley, \emph{Catalan Numbers}, Cambridge University Press, 2015. Two hundred and fourteen
interpretations.
\item W.~Feller, \emph{An Introduction to Probability Theory and Its Applications}, Vol.~1, 3rd ed., Wiley, 1968,
Chapters~XI--XII. Generating functions, random sums and branching.
\item D.~Aldous, \emph{The continuum random tree I}, Annals of Probability 19 (1991), 1--28. Large critical
Galton--Watson trees, seen from far away, converge to a random fractal tree.
\end{itemize}
